% ---------- sections/02_informe_economico.tex ----------
\chapter{Informe económico}
\label{ch:informe_economico}

\section{Resumen y enfoque}
\label{sec:resumen_enfoque}

Este capítulo presenta la situación económica al \textbf{31 de octubre de 2025} para el periodo \textbf{octubre 2023–octubre 2025}. El propósito es dejar un cierre claro y verificable, con cifras resumidas aquí y \textbf{soportes completos en anexos} (conciliaciones, extractos, comprobantes y actas). La lectura debe servir a la comunidad y a la vez a la auditoría interna y a la transición ordenada

\subsection*{Indicadores clave al corte}
\begin{table}[H]
\centering
\caption{Indicadores resumidos al 31/10/2025}
\label{tab:indicadores_resumen}
\begin{singlespace}
\begin{tabular}{p{9cm} r}
\toprule
\textbf{Concepto} & \textbf{Valor (USD)} \\
\midrule
Ingresos acumulados (Oct 2023–Oct 2025) & 184\,520.00 \\
Egresos acumulados (Oct 2023–Oct 2025) & 171\,845.60 \\
Saldo disponible al 31/10/2025 & 12\,674.40 \\
\midrule
Porcentaje de morosidad al corte & 6.8\% \\
Gasto operativo mensual promedio & 7\,235.25 \\
\bottomrule
\end{tabular}
\end{singlespace}
\end{table}

\subsection*{Notas de lectura}
\begin{itemize}
  \item Las cifras provienen de libros contables y extractos bancarios conciliados (ver Sección \ref{sec:conciliaciones})
  \item Los soportes firmados y certificados bancarios se adjuntan en anexos y fueron revisados con la comisión económica
\end{itemize}

% -------------------------------
\section{Conciliaciones bancarias (Nov 2024–Oct 2025)}
\label{sec:conciliaciones}

Se incluye el resumen de las conciliaciones de los últimos 12 meses, con correspondencia entre saldo contable y saldo bancario. Las diferencias responden a partidas identificadas y su regularización se documenta en la Sección \ref{sec:anexos_referencias}. El detalle mensual con firmas y certificaciones va en anexos

\subsection*{Resumen mensual (formato de ejemplo)}
\begin{table}[H]
\centering
\caption{Conciliaciones resumidas (formato de ejemplo)}
\label{tab:conciliaciones_resumen}
\begin{singlespace}
\begin{tabular}{p{2.4cm} r r r p{6.5cm}}
\toprule
\textbf{Mes} & \textbf{Saldo libros} & \textbf{Saldo banco} & \textbf{Diferencia} & \textbf{Partidas principales} \\
\midrule
Sep 2025 & 14\,215.80 & 13\,768.30 & 447.50 & Cheque en tránsito; depósito de fin de mes no acreditado \\
Oct 2025 & 12\,674.40 & 12\,352.10 & 322.30 & Nota de débito bancaria; ajuste de intereses \\
\bottomrule
\end{tabular}
\end{singlespace}
\end{table}

\subsection*{Checklist de soporte}
\begin{itemize}
  \item Conciliaciones mensuales firmadas y certificados de saldo
  \item Carpeta digital: \emph{anexos/C\_listados/conciliaciones\_2024\_2025/}
\end{itemize}

% -------------------------------
\section{Acta de arqueo de caja (31/10/2025)}
\label{sec:arqueo_caja}

Verificación física del efectivo en caja y valores en poder del tesorero al cierre del periodo, contrastando con los registros del libro diario

\subsection*{Cuadro de arqueo}
\begin{table}[H]
\centering
\caption{Arqueo de caja al 31/10/2025}
\label{tab:arqueo_caja}
\begin{singlespace}
\begin{tabular}{p{9cm} r}
\toprule
\textbf{Detalle} & \textbf{Valor (USD)} \\
\midrule
Saldo según libros al 31/10/2025 & 2\,450.00 \\
Efectivo contado & 2\,430.00 \\
\midrule
\textbf{Diferencia} & -20.00 \\
\bottomrule
\end{tabular}
\end{singlespace}
\end{table}

\subsection*{Observaciones y acciones}
\begin{itemize}
  \item Se detectó diferencia de USD 20.00 por reposición en tránsito que se registrará en noviembre 2025
  \item Adjuntar recibos pendientes y su estatus de registro según carpeta \emph{anexos/C\_listados/caja\_soportes/}
\end{itemize}

% \subsection*{Firmas}
% \vspace{1.5em}
% \noindent\rule{6cm}{0.4pt}\quad \rule{6cm}{0.4pt}\quad \rule{6cm}{0.4pt}
% 
% \noindent Tesorero \hfill Presidente \hfill Comisión de vigilancia

% -------------------------------
\section{Cuentas por cobrar (morosidad al corte)}
\label{sec:cuentas_por_cobrar}

Se presenta la cartera por antigüedad con fines de gestión. El \textbf{listado completo por usuario} va en anexos para proteger datos personales y permitir seguimiento caso por caso

\subsection*{Antigüedad de saldos}
\begin{table}[H]
\centering
\caption{Antigüedad de saldos de cartera}
\label{tab:cartera_antiguedad}
\begin{singlespace}
\begin{tabular}{p{9cm} r}
\toprule
\textbf{Rango} & \textbf{Monto (USD)} \\
\midrule
0–30 días & 6\,280.00 \\
31–60 días & 2\,145.00 \\
61–90 días & 1\,380.00 \\
91+ días & 3\,065.00 \\
\midrule
\textbf{Total cartera} & 12\,870.00 \\
\bottomrule
\end{tabular}
\end{singlespace}
\end{table}

\subsection*{Acciones sugeridas}
\begin{itemize}
  \item Acuerdos de pago escalonados para deudas mayores a 90 días
  \item Recordatorios y cortes programados conforme a reglamento y actas
\end{itemize}

\subsection*{Reportes de respaldo}
\begin{itemize}
  \item \emph{anexos/C\_listados/cartera\_al\_corte.xlsx} con detalle por usuario
  \item Copias de notificaciones y convenios firmados, si existen
\end{itemize}

% -------------------------------
\section{Cuentas por pagar (proveedores y servicios)}
\label{sec:cuentas_por_pagar}

Relación de obligaciones vigentes a la fecha de corte. El detalle de facturas y contratos se mantiene en anexos, incluyendo cronograma de pagos acordado con los proveedores

\subsection*{Resumen de obligaciones}
\begin{table}[H]
\centering
\caption{Cuentas por pagar al corte (ejemplo de formato)}
\label{tab:cxp_resumen}
\begin{singlespace}
\begin{tabular}{p{3.2cm} p{6cm} r p{2.8cm} p{2.4cm}}
\toprule
\textbf{Proveedor} & \textbf{Concepto} & \textbf{Monto (USD)} & \textbf{Vencimiento} & \textbf{Estado} \\
\midrule
Proveedor A & Insumos de cloración (ejemplo) & 1\,280.50 & 15/11/2025 & Pendiente/Parcial \\
Empresa B & Mantenimiento puntual (ejemplo) & 2\,450.00 & 05/12/2025 & Pendiente/Programado \\
\bottomrule
\end{tabular}
\end{singlespace}
\end{table}

\subsection*{Notas}
\begin{itemize}
  \item Adjuntar facturas, órdenes de compra o contratos como respaldo
  \item Carpeta digital sugerida: \emph{anexos/C\_listados/cxp\_soportes/}
\end{itemize}

% -------------------------------
\section{Notas metodológicas y supuestos}
\label{sec:notas_metodologicas}

Las cifras se elaboran con base en registros contables, extractos bancarios y comprobantes. Se prioriza la \textbf{trazabilidad} de cada total y la \textbf{comparabilidad} entre periodos. Cualquier ajuste posterior se documentará en acta y se reflejará en las siguientes conciliaciones

% -------------------------------
\section{Anexos y referencias internas}
\label{sec:anexos_referencias}

\begin{itemize}
  \item Conciliaciones y certificados: \emph{anexos/C\_listados/conciliaciones\_2024\_2025/}
  \item Cartera por usuario y soportes: \emph{anexos/C\_listados/cartera\_al\_corte.xlsx}
  \item Facturas/contratos de CxP: \emph{anexos/C\_listados/cxp\_soportes/}
\end{itemize}
% ---------- fin 02_informe_economico ----------
